% ---------------------------------------------------------------------------
% Academic CV -- one composition of the grammar.
%
% Everything here is a SLOT, not a claim. Names, institutions, venues and
% numbers are bracketed placeholders, and the bullets say what belongs in them,
% so the file shows how the primitives compose without asserting anything about
% a real person or organisation. Read it for shape; write content in main.tex.
%
% Compare examples/industry.tex: a different section list, the same seven
% primitives, the same shared right-hand edge.
%
% Build from the REPOSITORY ROOT so the class and fonts/ resolve:
%   latexmk -xelatex -outdir=build examples/academic.tex
% ---------------------------------------------------------------------------
\documentclass[fontset=commercial,pagefooter=false]{yuan2resume}

\cvname{[Full Name]}
\cvposition{[Ph.D. Candidate]}
\cvcontactrow{
  \cvcontact[\cvIconPhone]{[+00 0000 0000]}
  \cvemail{name@example.edu}
}
\cvcontactrow{
  \cvlink[\cvIconScholar]{https://example.edu/scholar}{Scholar}
  \cvlink[\cvIconGitHub]{https://example.com/git}{git/handle}
  \cvlink[\cvIconSite]{https://example.org}{example.org}
}

\begin{document}
\maketitle

\begin{cvtwocolumn}

\begin{cvsection}{Education}
  \cvpair{\textbf{[Doctoral Institution]}}{[City, Country]}
  \cvpair{\textit{Ph.D. Candidate in [Field]}}{Sept. 2020 -- \textit{Present}}
  \begin{cvitemize}
    \item Advisor and research area, one clause each -- this is the line a reader
      uses to place you, so it goes first.
    \item A second line only when it earns the space: a distinction, a
      fellowship, a thesis title.
  \end{cvitemize}
  \cvgap
  \cvpair{\textbf{[Undergraduate Institution]}}{[City, Country]}
  \cvpair{\textit{B.S. in [Field]}}{Sept. 2016 -- June 2020}
  \begin{cvitemize}
    \item One line at most. An older degree earns less space than a recent one.
  \end{cvitemize}
\end{cvsection}

\begin{cvsection}{Publications}
  \begin{cvnumlist}
    \item \textbf{[Your Name]}, {[Co-author]}. ``[Title of the First Paper].''
      \textit{[Venue]}, 2023.
    \item \textbf{[Your Name]}, {[Co-author]}, {[Co-author]}. ``[Title of the
      Second Paper, Long Enough That It Wraps to a Second Line].''
      \textit{[Venue]}, 2022.
    \item {[Co-author]}, \textbf{[Your Name]}. ``[Title of the Third Paper].''
      \textit{[Venue]}, 2021.
    \item \textbf{[Your Name]}, {[Co-author]}. ``[Title of the Fourth Paper].''
      \textit{[Venue]}, 2020.
  \end{cvnumlist}
\end{cvsection}

\begin{cvsection}{Research\\Experience}
  \cvpair{\textbf{[Lab, Institute or Company]}}{[City, Country]}
  \cvpair{\textbf{[Role]}, [Group or Team]}{2022 -- \textit{Present}}
  \begin{cvitemize}
    \item What you owned, in one line, with the scope made legible.
  \end{cvitemize}
  \begin{cventry}
    \cvpair{\textit{[Project or Workstream]}}{2023 -- 2024}
    \begin{cvitemize}
      \item The result, and the number that makes it checkable.
      \item A second line only when the project needs one. Note that the date
        above still sits on the same right edge as every other date.
    \end{cvitemize}
  \end{cventry}
\end{cvsection}

\begin{cvsection}{Skills}
  \textbf{Languages}: {\CodeFont{[Lang]}}, {\CodeFont{[Lang]}}, {\CodeFont{[Lang]}} \\
  \textbf{Frameworks}: [Framework], [Framework], [Framework], [Framework] \\
  \textbf{Spoken}: [Language] (native), [Language] (fluent), [Language] (basic)
\end{cvsection}

\begin{cvsection}{Awards\\and\\Honors}
  \cvpair{[Award or Fellowship Name]}{2022}
  \cvpair{[Award or Fellowship Name], [Awarding Body]}{2021}
  \cvpair{[Award or Fellowship Name]}{2020}
\end{cvsection}

\end{cvtwocolumn}
\end{document}
